% PAGES: 62-63  (printed 46-47)

\begin{table}[H]
\centering
\caption{Backward Substitution for upper triangular bidiagonal matrices}
\fbox{%
\begin{minipage}{0.82\textwidth}
\textbf{Function Call:}\\
$x = \text{backward\_subst\_bidiag}(U,b)$

\medskip
\textbf{Input:}\\
$U = $ nonsingular upper triangular bidiagonal matrix of size $n$\\
$b = $ column vector of size $n$

\medskip
\textbf{Output:}\\
$x = $ solution to $Ux = b$

\medskip
$x(n) = \dfrac{b(n)}{U(n,n)}$;\\
for $j = (n-1):1$\\
\hspace*{1.5em}$x(j) = \dfrac{b(j)-U(j,j+1)x(j+1)}{U(j,j)}$;\\
end
\end{minipage}%
}
\end{table}

\section{LU decomposition without pivoting}

The LU decomposition of a matrix provides a computationally efficient
way for solving linear systems; see section 2.5 for details. While the
LU decomposition without pivoting does not exist for every nonsingular
matrix, it is often used in practice, e.g., for tridiagonal matrices;
see section 2.5.1. Moreover, the LU decomposition without pivoting
contains the main idea for the recursive algorithm behind both the LU
decomposition with row pivoting, which exists for all nonsingular
matrices,\footnote{Note that, throughout this book, ``LU decomposition''
will refer to the LU decomposition without pivoting, unless otherwise
specified.} and the Cholesky decomposition; see section 2.6 and
section 6.1, respectively.

\begin{definition}
The $LU$ decomposition without pivoting of a nonsingular square matrix
$A$ consists of finding a lower triangular matrix $L$ with all entries
on the main diagonal equal to $1$ and a nonsingular upper triangular
matrix $U$ such that
\[
A \;=\; LU.
\]
The $L$ and $U$ matrices are called the $LU$ factors of $A$.
\end{definition}

Note that requiring the diagonal entries of the matrix $L$ to be equal
to $1$ is necessary for the uniqueness of the LU decomposition: The
$n\times n$ matrix $A$ has $n^2$ entries, while the matrix $L$ has
$\frac{n(n+1)}{2}$ entries below the main diagonal and on the main
diagonal, and the matrix $U$ has $\frac{n(n+1)}{2}$ entries above the
main diagonal and on the main diagonal, for a total of
$\frac{n(n+1)}{2}+\frac{n(n+1)}{2} = n^2+n$ entries. These entries can be
thought of as unknowns when looking for the LU decomposition of $A$.
Thus, we have $n^2+n$ unknowns to be found given the $n^2$ entries of
$A$. Requiring the $n$ entries of $L$ on the main diagonal to be equal
to $1$ reduces the number of unknowns to $n^2$, which creates the
framework for the LU decomposition to be unique; see Theorem 2.2.

\begin{definition}
Let $A$ be an $n\times n$ matrix. The leading principal minors of $A$
are the determinants of the $i\times i$ matrices $A_i = A(1:i,1:i)$ made
of the $i^2$ upper left entries of $A$, for $1\le i\le
n$.\footnote{For example, the leading principal minors of the matrix
$\begin{pmatrix} 2 & -3 & 0\\ 1 & 1 & 1\\ -1 & 5 & -3\end{pmatrix}$ are
\[
\det(2) = 2; \quad
\det\begin{pmatrix} 2 & -3\\ 1 & 1\end{pmatrix} = 5; \quad
\det\begin{pmatrix} 2 & -3 & 0\\ 1 & 1 & 1\\ -1 & 5 & -3\end{pmatrix} = -22;
\]
see section 5.2.1 for more details and examples.}
\end{definition}

\begin{theorem}
A matrix has an LU decomposition if and only if all the leading
principal minors of the matrix are nonzero.
\end{theorem}

\begin{proof}
The proof of this result is of no further relevance herein and can be
found, e.g., in Datta~\cite{12}.
\end{proof}

For the LU decomposition to be consistent, it should be unique, and this
is, indeed, the case:

\begin{theorem}
If it exists, the LU decomposition without row pivoting of a matrix is
unique.
\end{theorem}

The proof of this result can be found in the technical appendix from
Chapter 10; see the proof of Theorem 10.10 from section 10.6.

\subsection{Pseudocode and operation count for LU decomposition}

The $L$ and $U$ matrices from the LU decomposition of a matrix are
computed using a recursive algorithm, one row of $U$ and one column of
$L$ at a time.

Let $A$ be an $n\times n$ nonsingular matrix with LU decomposition
without pivoting. We are looking for an $n\times n$ lower triangular
matrix $L$ with main diagonal entries equal to $1$ and for an $n\times
n$ upper triangular matrix $U$ such that
\begin{equation}
LU \;=\; A.
\end{equation}

\begin{examplex}
\textit{$4\times 4$ matrix}\par
For clarity, we include explicit formulations of certain formulas for
$4\times 4$ matrices. In particular, formula (2.21) is written as
follows if $A$, $L$, and $U$ are $4\times 4$ matrices:
\[
\begin{aligned}
&\begin{pmatrix}
1 & 0 & 0 & 0\\
L(2,1) & 1 & 0 & 0\\
L(3,1) & L(3,2) & 1 & 0\\
L(4,1) & L(4,2) & L(4,3) & 1
\end{pmatrix}
\begin{pmatrix}
U(1,1) & U(1,2) & U(1,3) & U(1,4)\\
0 & U(2,2) & U(2,3) & U(2,4)\\
0 & 0 & U(3,3) & U(3,4)\\
0 & 0 & 0 & U(4,4)
\end{pmatrix}\\[4pt]
&\;=\;
\begin{pmatrix}
A(1,1) & A(1,2) & A(1,3) & A(1,4)\\
A(2,1) & A(2,2) & A(2,3) & A(2,4)\\
A(3,1) & A(3,2) & A(3,3) & A(3,4)\\
A(4,1) & A(4,2) & A(4,3) & A(4,4)
\end{pmatrix},
\end{aligned}
\]
since all the main diagonal entries of $L$ are required to be equal to
$1$.
\end{examplex}
